\chapter{HiPPO And Memory}
\label{ch:hippo-and-memory}

\section{Why HiPPO Appears In SSM Papers}

S4 did not begin from a random recurrent matrix. A central part of its success came from
initializing the state dynamics with matrices derived from HiPPO: High-order Polynomial
Projection Operators \citep{gu2020hippo}. HiPPO is a theory of online memory. It asks:

\[
  \text{How can a finite-dimensional state store useful information about an entire past signal?}
\]

This is exactly the question long-sequence models face. A neural time series may contain
relevant context spread across hundreds or thousands of time steps. A state vector cannot
store the raw past exactly unless it grows with sequence length, so it must store a
compressed summary. HiPPO gives a principled way to define such summaries.

The goal of this chapter is not to reproduce every HiPPO derivation. The original paper
develops a general framework involving measures, orthogonal polynomials, and projection
dynamics. For understanding S4 and S5, the most important idea is simpler:

\[
  \text{HiPPO initializes a recurrent state to track polynomial coefficients of the past.}
\]

Those coefficients form a compact approximation of the input history. The resulting
state matrices have a spectrum of timescales that is much better suited to long memory
than a generic random matrix.

\section{The Online Memory Problem}

Suppose a scalar signal $u(t)$ arrives over time. At time $t$, the full history is the
function
\[
  u_{\leq t}: \tau \mapsto u(\tau), \qquad 0 \leq \tau \leq t.
\]
Storing this full function is impossible with a fixed-size vector. Instead, we maintain a
state
\[
  \vx(t) \in \R^N
\]
that summarizes the history. The update should be online: when a new input arrives, we
update $\vx(t)$ without replaying the whole past.

This is the same abstract problem solved by an RNN hidden state. The difference is that
HiPPO specifies what the state should mean. Rather than letting a random recurrent matrix
learn memory from scratch, HiPPO defines $\vx(t)$ as coefficients of a projection of the
past signal onto a chosen basis.

The compression problem can be written informally as
\[
  u_{\leq t}(\tau)
  \approx
  \sum_{n=0}^{N-1} x_n(t)\,p_n(\tau;t),
\]
where $p_n$ are basis functions over the past interval and $x_n(t)$ are the coefficients
stored in the state.

\section{Projection Onto A Basis}

A basis lets us approximate a function by a finite list of coefficients. In ordinary
Euclidean geometry, projecting a vector onto an orthonormal basis gives coordinates:
\[
  \vv \approx \sum_{n=0}^{N-1} c_n \ve_n.
\]
For functions, the same idea uses inner products:
\[
  c_n = \langle f, p_n\rangle.
\]
If the basis functions $p_n$ are orthonormal under the chosen inner product, then the
coefficients $c_n$ give the best approximation to $f$ in the span of the first $N$ basis
functions.

HiPPO applies this idea to the input history. At time $t$, the state stores coefficients
of a projection of the past signal. The basis is usually built from orthogonal
polynomials, especially Legendre polynomials.

\section{A Quick Reminder On Orthogonality}

Two vectors are orthogonal if their dot product is zero. Two functions are orthogonal if
their inner product is zero:
\[
  \langle f,g\rangle
  =
  \int f(\tau)g(\tau)w(\tau)\,d\tau.
\]
The weight function $w(\tau)$ and interval of integration define the geometry. Different
choices produce different families of orthogonal polynomials.

Legendre polynomials are orthogonal on an interval with a uniform weight. The first few
on $[-1,1]$ are
\[
  P_0(x)=1,\qquad
  P_1(x)=x,\qquad
  P_2(x)=\frac{1}{2}(3x^2-1).
\]
Higher-order polynomials capture finer variation. If we project a signal onto the first
$N$ Legendre polynomials, low-order coefficients capture coarse trends and high-order
coefficients capture more detailed shape.

This gives a useful mental picture of HiPPO memory:
\[
  \vx(t)
  =
  \text{coefficients of a polynomial approximation to recent history}.
\]

\section{Why The Basis Changes Over Time}

The online memory problem is harder than a fixed projection problem because the interval
of history changes as time advances. At time $t$, the past is $[0,t]$. At time $t+\Delta$,
the past is $[0,t+\Delta]$. If we care more about recent history, the effective window or
weighting may also shift over time.

HiPPO handles this by deriving differential equations for how the projection coefficients
should change as new data arrive and as the underlying measure over history evolves.
Those coefficient dynamics have the state space form
\[
  \frac{d}{dt}\vx(t) = \mA(t)\vx(t)+\mB(t)u(t),
\]
with particular matrices determined by the polynomial family and measure.

For the most important variants used in SSM initialization, this can be turned into
time-invariant or nearly time-invariant dynamics after a change of variables. That is
what makes the connection to S4 possible.

\section{Legendre Memory Intuition}

The Legendre version of HiPPO is easiest to understand geometrically. Imagine always
keeping a polynomial approximation to the recent past. The constant coefficient stores
something like the average level. The linear coefficient stores something like the trend.
Quadratic and higher coefficients store curvature and finer temporal structure.

As a new input arrives, all coefficients must update:

\begin{itemize}
  \item the newest value must be inserted into the memory;
  \item older information must be shifted or rescaled because the time interval has
  changed;
  \item the polynomial approximation must remain the best projection under the chosen
  measure.
\end{itemize}

This produces a structured recurrent update. It is not arbitrary recurrence. It is the
update required to keep the polynomial approximation correct as time moves forward.

\section{Why This Helps Long Memory}

Random recurrent matrices do not automatically preserve useful information over long
horizons. They may forget too quickly, explode, or represent the past in a poorly
conditioned way. HiPPO gives a deliberately constructed memory system.

The benefit is not that polynomial projection is the only good way to remember. The
benefit is that it gives the model a strong starting point:
\[
  \text{state dimensions correspond to organized timescales and basis components}.
\]
When used as an initialization for a learnable SSM, the model does not have to discover
long memory from scratch. It starts from dynamics that already know how to compress
history.

For long-range sequence modeling, this matters. A model that begins with only short
timescale modes may never learn long dependencies well. HiPPO-derived matrices provide
a spectrum of modes that can represent both coarse and fine structure over long contexts.

\section{HiPPO As An SSM Initialization}

The key output of HiPPO for S4 is a state matrix $\mA$ and input matrix $\mB$. These
matrices define continuous-time dynamics:
\[
  \frac{d}{dt}\vx(t)=\mA\vx(t)+\mB u(t).
\]
The state $\vx(t)$ has an interpretation as projection coefficients. Then, as in
Chapter~\ref{ch:state-space-models}, the continuous system can be discretized and used
as a sequence layer.

In a deep SSM, the HiPPO matrices are not necessarily frozen. They are usually used to
initialize learnable parameters. Training can adapt them to the task. This is analogous
to initializing a convolutional filter bank with a useful signal-processing transform,
then allowing gradient descent to tune it.

S4 uses a particular HiPPO-inspired matrix as the starting point, then applies additional
structure so the resulting long convolution kernel can be computed efficiently. S5 keeps
the broad idea of S4-style initialization but simplifies the state parameterization.

\section{The HiPPO Matrix Shape}

The exact entries of the HiPPO matrix are not essential on a first pass, but it is useful
to know what kind of object it is. For the Legendre translated variant often discussed
around S4, $\mA$ is a structured lower-triangular matrix with entries determined by the
polynomial orders. It is not a generic learned dense matrix.

At a high level:

\begin{itemize}
  \item lower-order state coordinates track coarse components of the past;
  \item higher-order coordinates track finer components;
  \item the matrix couples these coordinates in the way required by online projection;
  \item the resulting dynamics have stable memory behavior across many timescales.
\end{itemize}

The original HiPPO paper defines several variants. They differ in the measure used to
weight the past. Some variants emphasize the entire history; others emphasize recent
history. For S4, the important fact is that HiPPO supplies a principled state matrix
whose memory properties are far better than an arbitrary initialization.

\section{Measures And Weighting The Past}

The word ``measure'' in HiPPO can sound abstract, but the modeling idea is intuitive.
When summarizing the past, we must decide which parts matter most. A uniform measure over
the whole past treats all times equally. A recency-weighted measure gives more importance
to recent inputs. A sliding-window-like measure focuses on a finite recent interval.

Different measures lead to different projection dynamics. This affects what the memory
state represents:
\[
\begin{array}{lll}
\text{whole-history weighting} & \rightarrow & \text{global summary of all previous input}, \\
\text{recency weighting} & \rightarrow & \text{stronger memory of recent input}, \\
\text{windowed weighting} & \rightarrow & \text{compressed representation of a recent interval}.
\end{array}
\]

For deep learning, the exact measure is less important than the fact that HiPPO turns a
choice about memory into concrete recurrent matrices.

\section{What HiPPO Does Not Solve}

HiPPO is not a full sequence model by itself. It solves an initialization and memory
representation problem. Several other problems remain:

\begin{itemize}
  \item How do we compute the resulting long convolution kernel efficiently?
  \item How do we parameterize the state matrix so it remains stable during training?
  \item How do we mix many input and output channels?
  \item How do we add nonlinear feature transformations around the linear memory?
  \item How do we scale the method on GPUs?
\end{itemize}

S4 answers these questions with structured state matrices and efficient kernel
computation. S5 answers them with a simpler diagonal/MIMO formulation and parallel scan.
Mamba changes the problem again by making the dynamics input-dependent.

\section{Connection To Neural Time Series}

For Utah array data, the signal is continuous in time but observed in bins. Speech motor
cortex activity has temporal structure across multiple scales: fast changes around
movement onset, slower context across phonemes or syllables, and session-level nuisance
variation. A useful encoder needs memory, but not just one memory timescale.

HiPPO-style initialization is attractive because it gives an SSM a structured set of
memory modes before task-specific training. In principle, some modes can track rapid
local changes while others retain slower context. Whether this helps a particular SSL
objective is empirical, but the architectural motivation is clear.

This also explains why SSMs are not merely ``RNNs with different branding.'' The point is
not just to use recurrence. The point is to use a recurrence whose dynamics are designed
to represent long histories and whose structure can be exploited computationally.

\section{A Minimal Worked Picture}

Suppose we approximate the recent past using three coefficients:
\[
  u_{\leq t}(\tau)
  \approx
  x_0(t)P_0(\tau)
  +
  x_1(t)P_1(\tau)
  +
  x_2(t)P_2(\tau).
\]
Very roughly:

\begin{itemize}
  \item $x_0(t)$ captures average level;
  \item $x_1(t)$ captures trend;
  \item $x_2(t)$ captures curvature.
\end{itemize}

When a new input arrives, the coefficients change. The update is not simply
``replace the oldest value.'' It must adjust the whole polynomial approximation. HiPPO
derives the continuous-time update that performs this coefficient maintenance online.

Increasing the state size $N$ adds more basis functions and therefore more detailed
memory. In an SSM layer, $N$ controls how rich this latent memory can be.

\section{What To Remember}

\begin{itemize}
  \item HiPPO is a framework for online compression of a signal history.
  \item The state stores coefficients of a projection of the past onto basis functions.
  \item Orthogonal polynomial bases, especially Legendre bases, give organized memory
  coordinates.
  \item HiPPO-derived matrices initialize SSMs with useful long-memory dynamics.
  \item The exact HiPPO derivation is less important at first than the memory principle:
  maintain a finite-dimensional approximation of the past as new inputs arrive.
  \item S4 builds on HiPPO but also needs additional structure for efficient long-kernel
  computation.
\end{itemize}

\section{Exercises}

\begin{enumerate}
  \item In your own words, explain the difference between storing the raw past and
  storing projection coefficients of the past.
  \item Why might a polynomial basis be useful for compressing a temporal signal?
  \item What kind of information would you expect low-order polynomial coefficients to
  store? What about high-order coefficients?
  \item Explain why a good recurrent initialization matters for long-context modeling.
  \item HiPPO gives a useful memory state, but not a complete deep sequence model. List
  two additional ingredients needed to turn it into an S4-style layer.
\end{enumerate}

\section{Prepares You For}

S4 uses structured versions of HiPPO-inspired state matrices to get both long memory and
efficient computation. The next chapter explains how S4 turns the SSM kernel
\[
  \mK_i=\mC\bar{\mA}^i\bar{\mB}
\]
into a practical deep-learning layer for long sequences.
